\section{Semantic Analysis}\label{sec:semantic-rules}

The semantic pass validates scoping, types, and musical constraints on the parser's AST.
It performs a single traversal: validate tempo and signature headers, register patterns when each scope opens,
then check statements and expressions, attaching symbol links and expression types that lowering reads without
re-analysis.
\topic{Scoping and writable boundaries.}
Symbols live in a stack of lexical environments (global, track, voice, pattern, control-flow block).
Lookup walks outward from the current scope for reads; writes are governed by a \textit{writable boundary} that
reflects whether an inner scope executes sequentially inside its parent or in parallel with it.

Sequential inner scopes inherit the parent's writable boundary.
Statements in a track body, patterns defined inside that track, and control-flow blocks such as \texttt{for},
\texttt{loop}, and \texttt{if} may therefore read and assign variables declared in the enclosing track (or
global scope when the boundary permits).
They execute in source order on the same timeline, so mutating a shared counter or motif parameter in the track
scope is well defined.

Parallel inner scopes receive a new writable boundary at their own scope.
A \texttt{voice} inside a track may \textit{read} track-level variables when building its part, but assignments
to those outer variables are rejected---the diagnostic reports a read-only outer binding.
Voices lower to parallel cursors on the same MIDI channel; allowing them to mutate the parent track's state would
conflate timelines that the language intentionally keeps independent.
Variables declared inside the voice remain fully writable within that voice.

Global patterns establish the writable boundary at pattern scope, so they cannot assign globals or track state
unless nested inside a track that has already set a wider boundary.
\topic{Types.}
Declarable types are \texttt{int}, \texttt{double}, \texttt{bool}, \texttt{note}, \texttt{seq}, and
\texttt{chord}.
Declarations and assignments require exact type equality: \texttt{int x = 1.0;} is rejected even though
arithmetic expressions may promote numeric operands.
Binary operators enforce numeric or boolean rules; sequences cannot be added with \texttt{+}.

\topic{Patterns and \texttt{play}.}
Pattern calls type-check arguments, then select the overload whose parameter types match.
\texttt{play} requires musical expression types on melodic tracks and drum literals on drum tracks; optional
\texttt{:duration} and \texttt{from} operands must be numeric.
These rules replace version~1.0's arity-only dispatch and second-pass pattern simulation.
